%  Workshop No. 4 - Interactive Dashboards with Plotly Dash
%  Volcanic Activity Monitoring System (United States), Team 6
%
%  Assembled by Part 2 (Julián Pulido) from the sections of the three parts.
%  Build: pdflatex report.tex (twice), from this folder. Figures are read from figures/.
%  Every section is written. The two Part 3 sections (storytelling dashboard, challenges) were drafted by
%  Part 2 from the Part 3 code and are marked '% Drafted by Part 2'; Part 3 reviews them.
%  \pending{owner}{source file} is kept for any section that has to be filled again.
\documentclass[11pt,a4paper]{article}

\usepackage[utf8]{inputenc}
\usepackage[T1]{fontenc}
\usepackage[english]{babel}
\usepackage{lmodern}
\usepackage{textcomp}
\usepackage[a4paper,margin=2.3cm]{geometry}
\usepackage{graphicx}
\usepackage{booktabs}
\usepackage{longtable}
\usepackage{array}
\usepackage{caption}
\usepackage{enumitem}
\usepackage{xcolor}
\usepackage{fancyhdr}
\usepackage{amsmath}
\usepackage{amssymb}
\usepackage[hidelinks]{hyperref}
\usepackage{tabularx}

\graphicspath{{figures/}{./}}

\hypersetup{
  pdftitle={Workshop 4 - Volcanic Activity Monitoring: Interactive Dashboards with Plotly Dash},
  pdfauthor={Team 6 - Universidad Distrital Francisco Jose de Caldas},
  pdfsubject={Data Analysis Programming 2026-II, Workshop No. 4}
}

\captionsetup{font=small,labelfont=bf}
\setlist[itemize]{itemsep=2pt,topsep=3pt}
\setlist[enumerate]{itemsep=2pt,topsep=3pt}

\pagestyle{fancy}
\fancyhf{}
\lhead{\small Workshop 4 --- Interactive Dashboards with Plotly Dash}
\rhead{\small Team 6}
\cfoot{\thepage}
\renewcommand{\headrulewidth}{0.4pt}

\newcolumntype{L}[1]{>{\raggedright\arraybackslash}p{#1}}

% Long file, column and test names in \texttt{} may break after an underscore, and paragraphs get a
% little extra stretch, so code names do not run into the margin
\renewcommand{\_}{\textunderscore\allowbreak}
\setlength{\emergencystretch}{3em}

% Finds images in figures/ or the root folder; renders a placeholder only if missing
\newcommand{\shot}[3]{%
  \begin{figure}[htbp]\centering
  \IfFileExists{figures/#1}{\includegraphics[width=#2\linewidth]{figures/#1}}{%
    \IfFileExists{#1}{\includegraphics[width=#2\linewidth]{#1}}{%
      \fbox{\parbox[c][3.2cm][c]{0.85\linewidth}{\centering\small Image pending: \texttt{\detokenize{#1}}}}}}
  \caption{#3}\label{fig:#1}
  \end{figure}}

% A section another part still has to deliver (removed when its text is pasted in)
\newcommand{\pending}[2]{%
  \begin{center}\fcolorbox{orange!70!black}{orange!8}{\parbox{0.9\linewidth}{\small
  \textbf{Pending} --- section owned by #1. Source: \texttt{\detokenize{#2}}.}}\end{center}}

\begin{document}

\begin{titlepage}
\centering
{\large Universidad Distrital Francisco José de Caldas\\ School of Engineering --- Computer Engineering\par}
\vspace{1.2cm}
{\Large Data Analysis Programming 2026-II\par}
\vspace{1.5cm}
{\huge\bfseries Workshop No.~4\\[4pt] Interactive Dashboards with Plotly Dash\par}
\vspace{0.8cm}
{\Large Volcanic Activity Monitoring System (United States)\par}
\vspace{1.5cm}
{\large Team 6\par}
\vspace{0.4cm}
\begin{tabular}{ll}
20231020006 & Juan Sebastián Rangel Pinto\\
20231020055 & Jhonatan David Vergara Jiménez\\
20231020064 & Julián David Pulido Carreño\\
\end{tabular}
\vfill
{\large Professor: Eng. Carlos Andrés Sierra, M.Sc.\par}
\vspace{0.3cm}
{\large Bogotá, \today\par}
\end{titlepage}

\tableofcontents
\clearpage

%==============================================================================
\section{Executive summary}\label{sec:summary}

Workshops 1--3 built a medallion pipeline for the volcanoes of the United States: Airflow DAGs land
the USGS (HANS alert notices, earthquake feed and FDSN history, Volcano Updates page) and Smithsonian GVP
(volcanoes, eruptions) data in Bronze, clean and deduplicate it in Silver, and PySpark writes two Gold
Parquet files, one for data governance and one for storytelling. Workshop~4 puts that Gold layer in front
of its two audiences with two Plotly Dash applications that read \emph{only} the newest Gold files:

\begin{itemize}
  \item \textbf{Governance dashboard} (\texttt{localhost:8050}, Section~\ref{sec:governance}) for the
        data team: nine quality KPIs against thresholds stored in \texttt{config/sources.yaml}, null and
        outlier rates per field, weekly volume per source and a data-health table per source.
  \item \textbf{Storytelling dashboard} (\texttt{localhost:8051}, Section~\ref{sec:storytelling}) for
        the emergency-management analyst: the activity tone of the USGS notices, the most frequent terms,
        API versus scraping, and the \emph{Volcano Activity Indicator} (VAI) panel.
\end{itemize}

The VAI (Section~\ref{sec:vai}) is the project's indicator: a 0--100 score per volcano that mixes the
official USGS alert level, the seismicity of the last seven days against the volcano's own weekly
history since 2010, and its eruptive state. On the 47 US eruptions since 2010 it was ELEVATED or higher
in the four weeks before the start of 52\,\% of them, with a false-alarm rate of 0.14 per quiet week
(0.06 for its four-week average). It is a monitoring aid, not an eruption forecast; the dashboard says so
next to it.

To feed the dashboards, Silver gained the end and the ``still going on'' flag of each eruption and the
USGS region of each volcano, and Gold reads Silver with a merged schema (Section~\ref{sec:integration}).
Both dashboards and the whole pipeline start with one \texttt{docker compose up}
(Section~\ref{sec:reproducibility}).

%==============================================================================
\section{Integration adjustments}\label{sec:integration}

The changes below were needed so the dashboards show accurate and complete data. They are kept apart
from the Workshop~3 pipeline improvements, as the workshop asks
(\texttt{workshop\_3/pipeline\_improvements.md}); the full notes are in
\texttt{workshop\_4/integration\_adjustments.md}.

\subsection{Silver eruptions: end date and continuing eruptions}
\textbf{Found:} the VAI gives an eruptive component of 1 to a volcano whose eruption is still going on,
but Silver kept only the start of each eruption. The GVP Bronze file has the end
(\texttt{EndDateYear/Month/Day}) and marks an eruption still going on at its last update with
\texttt{EndDateYearModifier = ">"} (for example Kilauea since 2024-12-23 and Great Sitkin since
2021-05-25).

\noindent\textbf{Change:} three new columns written by \texttt{eruptions\_from\_gvp}:
\begin{center}
\begin{tabular}{L{2.6cm}L{2.6cm}L{8.4cm}}
\toprule
\textbf{Column} & \textbf{Type} & \textbf{Meaning}\\
\midrule
\texttt{end\_time} & timestamp (UTC) & end of the eruption, first day of the known period; null when unknown\\
\texttt{end\_precision} & category & \texttt{day} / \texttt{month} / \texttt{year} / \texttt{unknown}, like \texttt{time\_precision}\\
\texttt{is\_continuing} & bool & true when GVP wrote \texttt{>} before the end year\\
\bottomrule
\end{tabular}
\end{center}
A missing end date alone does not mean ``ongoing'': 6,602 of the 11,089 GVP eruptions have none, almost
all of them old eruptions whose end is unknown. Only the \texttt{>} marker is read as continuing.

\subsection{Silver volcanoes: observatory and region}
\textbf{Found:} the VAI ranking must be filterable by region (user story US6: Alaska, Hawaii,
Cascades\ldots), but Silver \texttt{volcanoes} had no region.

\noindent\textbf{Change:} \texttt{subregion} (GVP, all 165 volcanoes), \texttt{observatory} and
\texttt{usgs\_region} (HANS volcano list, 153 of the 165). The dashboard shows the observatory group
(\texttt{indicator.regions} in \texttt{sources.yaml}), or the GVP subregion mapped to the same groups for
the 12 volcanoes HANS does not list. Every volcano gets one of six regions; none falls into ``Other''.

\subsection{Gold storytelling: four indicator sections}
\texttt{gold/indicator.py} adds \texttt{vai\_current}, \texttt{vai\_weekly}, \texttt{volcano\_map} and
\texttt{eruption\_history} to the storytelling file without changing the Gold contract. Two problems were
found while testing it:
\begin{itemize}
  \item GVP eruptions go back to year 1 and pandas cannot hold nanosecond timestamps before 1677, so
        collecting them from Spark failed with \texttt{OutOfBoundsDatetime}. Dropping old eruptions would
        change the result (a volcano that last erupted in 1510 has an eruptive component close to 0 but
        \emph{available}, which changes how the weights are shared), so the times now leave Spark as
        seconds since 1970.
  \item GVP day-precision dates fall on 00:00 UTC, and an eruption starting on a Monday was also counted
        in the week before (9 of the 47 US eruptions since 2010). Weeks are now half-open: Monday 00:00
        included, next Monday excluded.
\end{itemize}

\subsection{Gold reads Silver with \texttt{mergeSchema}}
\texttt{spark.read.parquet(folder)} takes the schema of one file. After a Silver column is added, a
folder holds files with and without it, and the new column could silently disappear from Gold.
\texttt{read\_topic} now reads each topic folder with \texttt{mergeSchema}.

\subsection{Applying them to an existing data lake}
Silver rewrites a record only when its source content changes, so the existing GVP rows are not rewritten
by later runs. On a data lake older than these changes the two GVP topics are rebuilt once with
\texttt{scripts/reprocess\_silver\_source.py gvp\_volcanoes} and \texttt{gvp\_eruptions} (with
\texttt{silver\_processing} paused), and then Gold is run again. A fresh clone needs nothing extra.

Each change has a regression test (\texttt{test\_eruption\_before\_1677\_still\_counts},
\texttt{test\_eruption\_starting\_on\_a\_monday\_does\_not\_reach\_the\_week\_before},
\texttt{test\_reader\_merges\_old\_and\_new\_silver\_schemas}, plus the Silver transform tests).

%==============================================================================
\section{Volcano Activity Indicator}\label{sec:vai}

The VAI is a 0--100 score per US volcano that says how active the volcano looks \emph{compared with its
own past}. It is a monitoring aid for the emergency-management analyst, \textbf{not an eruption
forecast}: official alert levels come from the USGS volcano observatories. Rules are in
\texttt{gold/vai\_rules.py}, the PySpark job in \texttt{gold/indicator.py}, every parameter in the
\texttt{indicator} block of \texttt{config/sources.yaml}, and the validation in
\texttt{notebooks/05\_vai\_validation.ipynb}.

\subsection{Formula}
\[
\mathrm{VAI} = 100\cdot\frac{w_R R + w_S S + w_H H}{\sum \text{weights of the available components}}
\]

\begin{longtable}{L{2.4cm}L{9.2cm}L{2.4cm}}
\toprule
\textbf{Component} & \textbf{Rule} & \textbf{Weight}\\
\midrule\endhead
$R$, reported activity & Latest USGS HANS alert level: NORMAL 0, ADVISORY 0.4, WATCH 0.7, WARNING 1.0.
  Not available when the volcano has no HANS notice. & 0.45\\
$S$, seismic anomaly & Earthquakes $M\ge 1.0$ within 30\,km in the last 7 days, compared with every complete
  week since 2010-01-04, \emph{weeks without earthquakes included}. Mid-rank percentile
  $p = (\text{weeks below} + \tfrac12\,\text{weeks equal})/\text{weeks}$; $S=\max(0,(p-0.5)/0.5)$.
  Not available below 20 earthquakes in the baseline. & 0.35\\
$H$, eruptive state & 1 while an eruption is going on (GVP \texttt{>} marker) or in a week an eruption
  overlaps; otherwise $e^{-\Delta t/25}$ with $\Delta t$ the years since the last eruption ended.
  Not available without a dated eruption. & 0.20\\
\bottomrule
\end{longtable}

\noindent Levels: LOW $<25\le$ MODERATE $<50\le$ ELEVATED $<75\le$ HIGH. A component that is not
available is left out and the other weights are shared again; the dashboard lists which components each
volcano used (US8).

\subsection{Decisions made while implementing it}
\begin{itemize}
  \item \textbf{Zero weeks in the baseline, mid-rank percentile.} Most volcanoes are silent most weeks.
        Without the zero weeks a volcano with no earthquake at all would get $S$ close to 1; with them,
        no earthquake now gives $p\le 0.5$ and $S=0$. The current week is not in the baseline because it
        is not over.
  \item \textbf{Weekly history from $S$ and $H$ only.} HANS keeps no past alert levels, so $R$ cannot be
        rebuilt for past weeks (US7).
  \item \textbf{Four-week average in the history chart.} One earthquake in a usually silent week already
        looks unusual; the average is easier to read and halves the false alarms.
  \item \textbf{No weekly-report bonus.} The optional GVP weekly RSS was never ingested, so the
        +0.2 / +0.1 bonus of the Workshop~1 design is not applied.
\end{itemize}

\subsection{What Gold holds}
\begin{center}
\begin{tabular}{L{3.0cm}L{3.4cm}L{7.2cm}}
\toprule
\textbf{Section} & \textbf{Rows (2026-11-25)} & \textbf{Content}\\
\midrule
\texttt{vai\_current} & 1,144 (143 volcanoes) & score, level, $R$, $S$, $H$ (null = not available), components used, earthquakes in the window, baseline percentile\\
\texttt{vai\_weekly} & 338,304 & score, $S$, $H$ per Monday since 2010 for 128 volcanoes\\
\texttt{volcano\_map} & 660 (165 volcanoes) & coordinates, current alert level and aviation colour code\\
\texttt{eruption\_history} & 918 & VEI and continuing flag of each eruption since 1900\\
\bottomrule
\end{tabular}
\end{center}
22 of the 165 volcanoes have no VAI: no HANS notice, fewer than 20 baseline earthquakes and no dated
eruption.

\subsection{Validation}
44 of the 47 US eruptions since 2010 happened at a volcano with a seismic baseline. Threshold: VAI
$\ge 50$ (ELEVATED). A false alarm is a week more than 8 weeks away from any eruption with VAI $\ge 50$.

\begin{center}
\begin{tabular}{lccc}
\toprule
\textbf{Signal} & \textbf{Hit ($\pm$4 weeks)} & \textbf{Early hit (4 weeks before)} & \textbf{False-alarm rate}\\
\midrule
Weekly VAI & 0.70 & 0.52 & 0.14\\
$S$ alone & 0.68 & 0.52 & 0.17\\
4-week average VAI & 0.50 & 0.36 & 0.06\\
\bottomrule
\end{tabular}
\end{center}

\begin{itemize}
  \item The Workshop~1 hit definition (before \emph{or after} the start) is partly circular, because $H=1$
        from the start week on. The early hit rate is the honest measure; at the default weights it equals
        $S$ alone, so the early signal comes from the earthquakes.
  \item \textbf{Case studies} (Figure~\ref{fig:vai_case_studies.png}): Kilauea (summit eruption of
        2020-12-20), Mauna Loa (2022-11-27, months of rising seismicity before) and Great Sitkin
        (2021-05-25) were ELEVATED or HIGH before the start. Ahyi (2025-10-20), a submarine volcano with
        sparse seismic coverage, was missed ($S=0$). Kilauea 2018 cannot be tested: GVP counts it as the
        end of the 1983--2018 eruption, which started before the baseline.
  \item \textbf{Weights} (Figure~\ref{fig:vai_weight_sensitivity.png}): giving $H$ more than half of the
        weight raises the early hit rate to 0.82--0.89, but only because frequently erupting volcanoes keep
        $H$ near 1 for years and sit at ELEVATED most of the time. An equal split of $S$ and $H$ gives
        fewer false alarms (0.12 against 0.14) with almost the same early hit rate (0.50 against 0.52).
        The Workshop~1 weights are kept; the equal split is the alternative to discuss. At the HIGH
        threshold the false-alarm rate stays near 0.05 for any weight.
\end{itemize}

\shot{vai_hits_false_alarms.png}{0.85}{Hit, early-hit and false-alarm rates of the VAI around the US eruptions since 2010.}
\shot{vai_case_studies.png}{0.95}{Weekly VAI and its four-week average around four US eruptions.}
\shot{vai_weight_sensitivity.png}{0.8}{Early hit and false-alarm rates when the seismic share of the weight changes.}

\subsection{Limitations}
\begin{itemize}
  \item The weekly $S$ is in-sample: each past week is compared with the whole 2010--present history.
        The current $S$ uses only past weeks.
  \item Earthquakes go to the nearest volcano within 30\,km; in clustered fields (Aleutians, Cascades) a
        swarm under one volcano can raise its neighbour too.
  \item Volcanoes with sparse networks (submarine, Northern Mariana Islands) rarely reach 20 baseline
        earthquakes, so their VAI rests on $R$ and $H$.
  \item With only one component available, renormalisation gives it the whole weight, so a few
        unmonitored Alaska volcanoes with one or two recent earthquakes top the ranking. The ranking draws
        those single-input scores lighter and its hover gives the number of inputs.
\end{itemize}

%==============================================================================
\section{Governance dashboard documentation}\label{sec:governance}

The governance dashboard (\texttt{dashboard/governance\_app.py}, \texttt{http://localhost:8050}) shows the
data quality of the pipeline to the \emph{data team}: the people who run the DAGs and must decide whether
the data is fit to feed the storytelling dashboard. It answers three questions in this order:
\emph{is everything within its threshold?} (KPI cards and summary table), \emph{where is the problem?}
(charts per field, numeric field and source) and \emph{is every source still arriving?} (data-health
table, user story US4 of Workshop~1).

It reads only the newest \texttt{datalake\_gold/governance\_*.parquet} with \texttt{pd.read\_parquet()}
through the shared loader of Part~1; every number is computed in the Gold layer with PySpark and the app
only filters, sorts and formats. It starts with \texttt{python -m dashboard.governance\_app} or as the
\texttt{governance-dashboard} service of \texttt{docker compose up}, and refreshes with the
\emph{Refresh data} button or every five minutes.

\shot{governance_dashboard.png}{0.92}{Governance dashboard at 1920\,px width with the Gold file of
2026-11-28 02:56 UTC: header, filters, KPI cards, null and outlier rates, volume per source, records per
topic, quality summary table and data health per source.}

\subsection{Layout}
The page is a Bootstrap grid (Dash Bootstrap Components, theme Flatly) that fits a 1920$\times$1080 screen
without horizontal scrolling. From top to bottom:
\begin{enumerate}
  \item \textbf{Header} (shared with the storytelling app): project title, dashboard name, topic angle,
        team and the \texttt{generated\_utc} time of the Gold file shown.
  \item \textbf{Controls}: refresh button, and two filters, \emph{Source} and \emph{Silver topic}.
  \item \textbf{Six KPI cards}: total records, null rate, duplicate rate, schema compliance, quarantine
        rate and freshness, each with its value, a \emph{Pass}/\emph{Warning} badge and its threshold.
  \item \textbf{Two charts per row}: null rate per field $|$ outlier rate per numeric field; records per
        week and source $|$ records stored per topic.
  \item \textbf{Two tables}: quality summary (every KPI against its threshold) $|$ data health per source.
\end{enumerate}

\subsection{Component list}
\begin{longtable}{L{3.7cm}L{2.2cm}L{2.6cm}L{5.9cm}}
\toprule
\textbf{Component (id)} & \textbf{Type} & \textbf{Gold section} & \textbf{What it answers / interaction}\\
\midrule\endhead
KPI cards (\texttt{kpi-cards}) & Bootstrap cards & \texttt{kpi\_summary} & Is each headline KPI within its threshold? Badge with text, not colour only\\
Null rate (\texttt{null-rate}) & Horizontal bars & \texttt{null\_rate} & Which fields are the emptiest? 20 worst fields, sorted by severity; \emph{topic} filter; hover gives the exact \%\\
Outlier rate (\texttt{outlier-rate}) & Grouped bars & \texttt{outlier\_rate} & Outliers per numeric field with the KPI criterion (per-volcano fences) next to Silver's per-file flags, and the threshold line\\
Volume (\texttt{volume-over-time}) & Lines, log axis & \texttt{volume\_weekly} & Records per week and source since 2010; \emph{source} filter; zoom and hover per week\\
Records per topic (\texttt{records-per-topic}) & Horizontal bars, log axis & \texttt{volume\_total} & How many records each Silver topic and the quarantine hold\\
Quality summary (\texttt{quality-summary}) & Table & \texttt{kpi\_summary} & All nine indicators: what is measured, value, threshold, status\\
Data health (\texttt{data-health}) & Table & \texttt{data\_health}, \texttt{freshness} & Per source: newest file, file age, watermark age, records in Silver, rejected records; \emph{source} filter\\
\bottomrule
\end{longtable}

\subsection{KPI mapping}
Each KPI of Workshop~3 is shown with one headline value, computed in Gold (section \texttt{kpi\_summary})
from the rows of the other indicators, so the cards, the table and the charts can never disagree. The
thresholds are configuration (\texttt{config/sources.yaml}, \texttt{governance.thresholds}) and each one
is grounded in the baseline measured in Workshop~3 (findings F01--F24).

\begin{longtable}{L{2.7cm}L{4.1cm}L{1.9cm}L{1.5cm}L{4.0cm}}
\toprule
\textbf{KPI} & \textbf{Headline value} & \textbf{Threshold} & \textbf{Value} & \textbf{Why this threshold}\\
\midrule\endhead
Null rate & mean \% of empty values per field & $\le 5\,\%$ & 4.1\,\% \checkmark & baseline 3.6\,\% from known structural gaps (F01--F04); 4.1\,\% with the Workshop~4 columns (see below)\\
Duplicate rate & \% exact copies in Silver & $\le 0\,\%$ & 0.089\,\% \textbf{!} & must be 0; the feed$\leftrightarrow$FDSN hash bug (F08) is still open\\
Schema compliance & lowest \% of files matching the contract, per topic & $\ge 100\,\%$ & 100\,\% \checkmark & every file must follow the contract (it caught a schema drift, see below)\\
Quarantine rate & \% of processed records rejected & $\le 2\,\%$ & 0.70\,\% \checkmark & baseline 0.68\,\%, all business rules\\
Outlier rate & highest \% of outliers of a field & $\le 15\,\%$ & 9.4\,\% \checkmark & baseline 9.4\,\% (depth, real slab events)\\
Encoding errors (text\_length) & \% of notices with garbled text & $\le 0.1\,\%$ & 0.054\,\% \checkmark & baseline 0.054\,\% (F21)\\
Freshness & \% of sources with a file younger than a day (2$\times$ schedule if longer) & $\ge 100\,\%$ & 100\,\% \checkmark & no source may stay silent for more than a day\\
Ingestion frequency & \% of the last 7 days' hours with a new \texttt{usgs\_feed} file & $\ge 80\,\%$ & 0\,\% \textbf{!} & the feed changes every hour\\
Total records & records stored in Silver & --- & 311,083 & context, no threshold\\
\bottomrule
\end{longtable}

\noindent The two warnings (\textbf{!}) are real findings, not defects of the dashboard:
\begin{itemize}
  \item \textbf{Duplicate rate}: the dedup index of Silver keeps only the last hash per earthquake, and the
        feed and FDSN hash the same event differently (Workshop~3, F07--F08).
  \item \textbf{Ingestion frequency}: the Docker stack was stopped from 21 to 27 November, so no feed file
        arrived in the last seven full days before the run.
\end{itemize}

\paragraph{The schema-compliance KPI at work.} On the first run of the dashboard (Gold of 2026-11-28
02:42 UTC) schema compliance was \textbf{0\,\%}: the Silver files of \texttt{volcanoes} and
\texttt{eruptions} had been written on 2026-11-17, before Workshop~4 added six columns to their contract
(\texttt{subregion}, \texttt{observatory}, \texttt{usgs\_region}, \texttt{end\_time},
\texttt{end\_precision}, \texttt{is\_continuing}; Section~\ref{sec:integration}). The GVP catalogue is
downloaded once a month, so those files were never rewritten. The two sources were reprocessed with
\texttt{scripts/reprocess\_silver\_source.py} (Silver DAG paused meanwhile) and the next Gold run gave
100\,\% in all seven topics. The new columns also raised the null rate from 3.6\,\% to 4.1\,\%, an expected
change: \texttt{eruptions.end\_time} is empty for 69.7\,\% of the eruptions, whose end GVP does not know,
and \texttt{volcanoes.observatory} for the 7.3\,\% of volcanoes outside the HANS list.

\subsection{Gold additions for the dashboard}
\begin{itemize}
  \item \texttt{kpi\_summary}: per KPI the headline value (the metric name says what it is, e.g.
        \texttt{mean\_pct\_null}), \texttt{threshold\_max} or \texttt{threshold\_min} and \texttt{status}
        (1 pass, 0 warning).
  \item \texttt{data\_health} (US4): per source the records in Silver, the rejected records, the time of the
        newest Bronze file and the watermark age, measured as the age of the newest source timestamp stored
        in Silver, so Gold still reads Silver only.
  \item \texttt{volume\_weekly} gained the metric \texttt{records\_by\_source}, because the workshop asks
        for volume per source and \texttt{earthquakes} mixes two sources.
  \item Performance: \texttt{kpi\_summary} reads the Gold rows about ten times; materialising them with
        \texttt{localCheckpoint()} cut the Spark test on the Silver samples from 142\,s to 22\,s.
\end{itemize}

\subsection{Design rationale}
\begin{itemize}
  \item \textbf{Overview first, then detail.} Cards and the summary table answer ``is anything wrong?''
        in one glance; the charts and the data-health table explain where.
  \item \textbf{One source of truth.} Status is computed in Gold, not in the browser, so the dashboard, the
        report figures and any other consumer read the same verdict.
  \item \textbf{Thresholds are configuration}, versioned with the code, each with its baseline in a comment.
  \item \textbf{Status never relies on colour alone}: every badge says \emph{Pass}, \emph{Warning} or
        \emph{Info}.
  \item \textbf{Robust to missing data}: with an empty \texttt{datalake\_gold/}, or a Gold file written
        before a section existed, every component shows a message instead of failing.
\end{itemize}

\subsection{Testing}
\texttt{tests/test\_governance\_dashboard.py} (11 tests) builds every component from a small Gold file with
known values: card values and badges, sorting and filters of the charts, axis titles and the source note,
the fallback for older Gold files, the empty-Gold case and the callback that fills the eight outputs.
\texttt{tests/test\_gold\_governance.py} adds three Spark tests for \texttt{data\_health},
\texttt{kpi\_summary} (including cases that must warn) and the thresholds in \texttt{sources.yaml}.
The full suite passes (216 tests).

%==============================================================================
\section{Storytelling dashboard documentation}\label{sec:storytelling}
% Drafted by Part 2 (report assembly) from dashboard/storytelling_app.py, gold/storytelling.py,
% tests/test_storytelling_dashboard.py and notebooks/07_storytelling_figures.ipynb (branch w4-storytelling,
% commit 70cc9e9). To be reviewed by Part 3 (Juan Sebastián).

The storytelling dashboard (\texttt{dashboard/storytelling\_app.py}, \texttt{http://localhost:8051})
answers one question for the emergency-management analyst of Workshop~1: \emph{which US volcanoes show
unusual activity, and is it building or fading?} It reads only the newest
\texttt{datalake\_gold/storytelling\_*.parquet}, in plain language and without pipeline jargon; the
technical view of the same data is the governance dashboard (Section~\ref{sec:governance}).

\shot{storytelling_dashboard.png}{0.85}{Storytelling dashboard at \texttt{localhost:8051}: narrative
card, tone distribution, filters, tone trajectory, weekly volume, top keywords, API vs.\ scraping and the
Volcano Activity Indicator panel.}

\subsection{Story flow and layout}
The page follows the story flow agreed in the plan: \emph{overview} $\rightarrow$ \emph{what changed this
week} $\rightarrow$ \emph{which volcanoes and why}.
\begin{enumerate}
  \item \textbf{Overview}: the \emph{Operational Narrative Summary} card, next to the distribution of the
        activity tone of all notices.
  \item \textbf{What changed}: three filters (tone scope: all volcanoes or one volcano with an elevated
        alert; volume source; single keywords or bigrams), then the weekly tone trajectory, the weekly
        volume per source, the most frequent terms coloured by tone and the comparison of the two text
        sources.
  \item \textbf{Which volcanoes and why}: the Volcano Activity Indicator panel of Part~1
        (Section~\ref{sec:vai}): map of the official alert levels, ranking by indicator filterable by
        region, and the history of one volcano with its eruptions.
\end{enumerate}

\subsection{Gold additions}
Three sections were added to \texttt{gold/storytelling.py} for the dashboard:
\begin{itemize}
  \item \texttt{tone\_distribution}: notices per tone label and their share. The label comes from the
        domain lexicon score of each notice summary (Workshop~3): escalating above $+0.15$, calming below
        $-0.15$, stable in between.
  \item \texttt{volume\_activity}: records per week and source (HANS API notices, scraped Volcano
        Updates, earthquakes near a volcano).
  \item \texttt{narrative\_highlights}: the few numbers the narrative card quotes: mean tone of the latest
        week and its change against the previous week, the volcano with the most escalating language that
        week, and the number of volcanoes with an elevated official alert.
\end{itemize}

\subsection{Component list and user-story mapping}
\begin{longtable}{L{3.8cm}L{2.2cm}L{2.9cm}L{1.5cm}L{3.9cm}}
\toprule
\textbf{Component (id)} & \textbf{Type} & \textbf{Gold section} & \textbf{Story} & \textbf{What it tells the analyst}\\
\midrule\endhead
Narrative card (\texttt{narrative-card}) & Text card with a badge & \texttt{narrative\_highlights}, \texttt{tone\_distribution}, \texttt{alerts\_breakdown} & US1, US7 & Elevated alerts now, whether the language of the latest week is building or fading (badge \emph{Unrest Building}, \emph{Steady Posture} or \emph{Unrest Fading}), the volcano to watch\\
Tone distribution (\texttt{tone-donut}) & Donut, three categories & \texttt{tone\_distribution} & context & Share of escalating, stable and calming notices since 2010\\
Tone trajectory (\texttt{tone-trend}) & Lines + threshold bands & \texttt{tone\_weekly} & US7 & Weekly tone of all volcanoes or of one volcano, lexicon next to the VADER baseline\\
Weekly volume (\texttt{volume-activity}) & Lines, peaks marked & \texttt{volume\_activity} & US2, US3 & Records per week and source, with the activity peaks annotated\\
Top keywords (\texttt{top-keywords}) & Horizontal bars & \texttt{top\_keywords} & context & The 15 most frequent terms or bigrams, coloured by their mean tone\\
API vs.\ scraping (\texttt{source-comparison}) & Grouped bars, two panels & \texttt{source\_comparison} & context & Text length and tone of HANS notices against the scraped Volcano Updates page\\
VAI panel (Part~1) & Map, ranked bars, history & \texttt{vai\_current}, \texttt{vai\_weekly}, \texttt{volcano\_map}, \texttt{eruption\_history} & US1, US3, US5, US6, US7, US8 & Which volcanoes are unusual compared with their own past, and since when\\
\bottomrule
\end{longtable}

\subsection{Narrative rationale}
\begin{itemize}
  \item \textbf{Words before charts.} The narrative card is generated from Gold on every refresh and says
        in four sentences what the charts below show, so an analyst with two minutes still gets the
        message.
  \item \textbf{One colour code for tone everywhere}: calming green, stable grey, escalating red
        (\texttt{TONE\_COLORS} of the shared theme), in the donut, the trajectory bands and the keyword bars.
  \item \textbf{The indicator is not a forecast}, and the panel says so above the map, next to the official
        USGS alert levels, which are shown unchanged.
  \item The donut is the only circular chart in the project, kept because it has exactly three categories
        that add up to a whole; every comparison with more categories uses bars.
\end{itemize}

\subsection{Testing}
\texttt{tests/test\_storytelling\_dashboard.py} checks the narrative card with real Gold rows and its
fallback without data, that every storytelling figure has a title and the source annotation, and that the
layout contains the filters and the cards. Static versions of the main charts are drawn with Seaborn in
\texttt{notebooks/07\_storytelling\_figures.ipynb} (Figures~\ref{fig:story_tone_trajectory_distribution.png}--\ref{fig:story_source_comparison_volume.png}).

\shot{story_tone_trajectory_distribution.png}{0.95}{Weekly activity tone (lexicon vs.\ VADER) and tone distribution of the notices.}
\shot{story_top_keywords_by_tone.png}{0.95}{Most frequent terms and bigrams, coloured by their mean activity tone.}
\shot{story_source_comparison_volume.png}{0.95}{HANS API vs.\ scraped Volcano Updates: text length, tone and weekly volume.}

%==============================================================================
\section{Visualization design decisions}\label{sec:design}

\subsection{Chart types}
\begin{longtable}{L{3.6cm}L{3.0cm}L{7.6cm}}
\toprule
\textbf{Question} & \textbf{Chart} & \textbf{Why}\\
\midrule\endhead
Which fields are the emptiest? & Horizontal bars, sorted & The answer is a ranking and the labels (\texttt{topic.column}) are long: horizontal bars keep them readable and the order shows severity\\
How do two criteria compare per field? & Grouped bars + threshold line & Side-by-side bars make the gap between the per-volcano KPI and Silver's per-file flags visible (F19)\\
How does volume change over 16 years? & Lines & A time series; one line per source\\
Several sources of very different size & Log axis, ticks at powers of ten & Earthquakes ($\sim$1,000/week) and notices ($\sim$20/week) differ by two orders of magnitude; a linear axis flattens the smaller source\\
Is each KPI within its threshold? & Cards + table & Exact values matter more than shape; a table holds all nine KPIs with units\\
How old is each source's data? (report) & Heatmap, neutral blues & Many sources $\times$ two ages at a glance; colour is not red because an age is only bad relative to its schedule\\
\bottomrule
\end{longtable}
\noindent No pie charts: every categorical comparison has more than three categories, and bars are read
more accurately than angles, as the workshop recommends.

\subsection{Colour}
Both dashboards and the report figures use one palette, defined once in
\texttt{dashboard/\allowbreak common/\allowbreak theme.py}: Okabe--Ito, designed to stay distinguishable for the common types of
colour blindness. Pass is green (\#009E73) and warning vermillion (\#D55E00), always with a text label.
The activity tone of the storytelling dashboard keeps the workshop rule: calming green, stable grey,
escalating red. Alert levels and aviation colour codes keep their USGS meaning.

\subsection{Layout and interaction}
\begin{itemize}
  \item A shared header and Plotly template make the two apps look like one product.
  \item Every chart has a descriptive title, axis titles with units and a ``Source:'' note
        (\texttt{style\_figure}); hover gives exact values, the toolbar zooms and resets.
  \item Filters only appear where the data has that breakdown, and a line under them says which charts
        they affect.
  \item A grid of at most two charts per row keeps the page readable at 1080p without horizontal scrolling.
\end{itemize}

\subsection{Decisions changed after reviewing the screenshots}
\begin{itemize}
  \item The volume chart first had a range slider; it covered the source note and pushed the axis title
        down, and the toolbar already zooms, so it was removed.
  \item The log axis first labelled 2, 5, 100, 2, 5\,\ldots; ticks at powers of ten only are easier to read.
  \item Freshness first required a new file within twice the schedule interval (2\,h for hourly sources)
        and warned all the time: Bronze writes a file only when the source changed (F23) and HANS does not
        publish every hour. A source is now fresh when its newest file is younger than a day, or than twice
        its schedule when that is longer; the hourly gaps of the feed are what ingestion frequency checks.
  \item The data-health heatmap of the report first used a red scale and painted the monthly GVP catalogue
        dark red at 265\,h, a false alarm; it now uses neutral blues and leaves the verdict to the KPI.
\end{itemize}

\subsection{Static figures for the report}
\texttt{notebooks/06\_governance\_figures.ipynb} draws Seaborn versions of the governance charts from the
same Gold file and palette. Figures~\ref{fig:gov_kpi_summary.png}--\ref{fig:gov_data_health.png} show
the main ones.

\shot{gov_kpi_summary.png}{0.95}{Governance KPIs against their thresholds (Gold of 2026-11-28 02:56 UTC).}
\shot{gov_null_rate.png}{0.7}{Share of empty values per field, coloured by Silver topic.}
\shot{gov_outlier_rate.png}{0.6}{Outlier rate per numeric field: per-volcano fences (KPI) vs.\ Silver flags.}
\shot{gov_volume_weekly.png}{0.95}{Records per week and source since 2010 (log scale).}
\shot{gov_duplicates.png}{0.75}{Duplicates before and after deduplication.}
\shot{gov_data_health.png}{0.95}{Age of each source's data, and records and rejects per source (US4).}

%==============================================================================
\section{Challenges and resolutions}\label{sec:challenges}
% Drafted by Part 2 (report assembly) from the work of the three parts; to be reviewed by Part 3.

\begin{longtable}{L{4.2cm}L{6.4cm}L{3.6cm}}
\toprule
\textbf{Challenge} & \textbf{Resolution} & \textbf{Status}\\
\midrule\endhead
Silver files written before the Workshop~4 contract (six new columns) & The schema-compliance KPI showed 0\,\%; the two monthly GVP sources were reprocessed from Bronze with \texttt{scripts/reprocess\_silver\_source.py} & Resolved (100\,\%)\\
Freshness warned all the time & Bronze writes only when a source changes and HANS does not publish every hour; a source is now fresh below one day (or twice its schedule) and the hourly gaps of the feed are checked by ingestion frequency & Resolved\\
Gold summary took minutes in Spark & The KPI summary re-read the plan of eleven indicators; \texttt{localCheckpoint()} materialises the rows once (142\,s $\rightarrow$ 22\,s in the test) & Resolved\\
Storytelling Spark tests failed outside Docker & Spark runs the tone UDFs in separate Python workers that did not see \texttt{src/}; \texttt{tests/conftest.py} now exports \texttt{PYTHONPATH} and \texttt{PYSPARK\_PYTHON} & Resolved\\
Java 25 on one laptop is not supported by Spark 3.5 & Spark tests skip themselves on such a host and run inside the Airflow image (Java 17) & Worked around\\
7.5\,GB of RAM for Airflow, two databases, Spark and the tests & The two Gold tasks run one after the other; long test runs get a time limit; unused services are stopped while testing & Worked around\\
Three parts building charts in parallel & The long Gold contract (one \texttt{section} per chart) let each part add sections without touching the others' code; the branches merged without conflicts & Resolved\\
Duplicate earthquakes in Silver (feed and FDSN hash the same event differently) & Detected by the duplicate-rate KPI (warning); Gold counts the newest version only, so the dashboards are not affected; fix proposed for Silver & Open (Part~1)\\
A week without ingestion (stack stopped 21--27 November) & Visible as an ingestion-frequency warning; documented instead of hidden & Clears after a week of normal runs\\
Tone labels lean to ``escalating'' (78.3\,\% of notices) & The $\pm0.15$ lexicon threshold treats ordinary monitoring vocabulary (\emph{lava}, \emph{eruption}, \emph{activity}) as escalation; the VADER baseline is shown next to it so the reader can compare & Open: calibrating the thresholds per percentile is the next step\\
\bottomrule
\end{longtable}

%==============================================================================
\section{Reproducibility evidence}\label{sec:reproducibility}

\subsection{Setup from a fresh clone}
The README (\emph{Getting started on your own machine}) lists every step; in short:
\begin{enumerate}
  \item Prerequisites: Git, Docker Desktop with at least 4\,GB of RAM (5\,GB recommended for the Gold
        job), Python 3.11 for the tests and notebooks. On Windows, a repository path under about 95
        characters or long paths enabled.
  \item \texttt{git clone https://github.com/jhonatanvj/Proyecto\_Analisis\_De\_Datos.git}, then
        \texttt{cp .env.example .env} and set \texttt{HTTP\_USER\_AGENT} (USGS and GVP ask clients to
        identify themselves). \texttt{.env} is git-ignored.
  \item \texttt{docker compose up -d --build} starts Airflow (\texttt{localhost:8080}), the PostGIS
        project database with its migrations, and both dashboards (\texttt{localhost:8050} and
        \texttt{8051}).
  \item In Airflow, unpause \texttt{bronze\_ingestion} (hourly) and \texttt{silver\_processing}
        (minute 15); run \texttt{bronze\_backfill\_fdsn} once for the earthquake history since 2010;
        trigger \texttt{gold\_processing} (weekly, Monday 03:00 UTC) or run it with
        \texttt{airflow dags test gold\_processing}.
  \item Press \emph{Refresh data} in either dashboard, or wait up to five minutes, to read the new Gold
        files.
\end{enumerate}

\subsection{Checks}
\begin{itemize}
  \item A fresh clone into a short path, with a new virtual environment, installs
        \texttt{requirements.txt} and passes the test suite. With the three parts merged, a host with
        Java 17 runs the whole suite: 219 passed, 1 skipped (2026-11-28). On a host with Java 25, which
        Spark 3.5 does not support, the Spark tests skip themselves and run inside the Airflow container,
        which has Java 17.
  \item Library versions are bounded below the next major release (\texttt{pandas<4}, \texttt{numpy<3},
        \texttt{dash<5}), and the suite passes with the newest pandas 3 release.
\end{itemize}

\shot{airflow_dags_w4.png}{0.95}{Airflow: the Bronze, Silver and Gold DAGs after the end-to-end run.}
\shot{airflow_gold_grid_w4.png}{0.95}{Airflow: grid of the \texttt{gold\_processing} run that wrote the Gold files shown in the dashboards.}
\shot{governance_dashboard.png}{0.75}{Governance dashboard at \texttt{localhost:8050}.}
\shot{storytelling_dashboard.png}{0.55}{Storytelling dashboard at \texttt{localhost:8051}.}
\shot{storytelling_vai_panel.png}{0.95}{Volcano Activity Indicator panel: map, ranking by region and history of one volcano.}

%==============================================================================
\section*{References}
\addcontentsline{toc}{section}{References}
\begin{itemize}
  \item Okabe, M. and Ito, K. (2008). \emph{Color Universal Design (CUD): How to make figures and
        presentations that are friendly to colorblind people}. https://jfly.uni-koeln.de/color/
  \item Shneiderman, B. (1996). The eyes have it: a task by data type taxonomy for information
        visualizations. \emph{Proc. IEEE Symposium on Visual Languages}, 336--343.
  \item Few, S. (2006). \emph{Information Dashboard Design}. O'Reilly.
  \item Plotly Technologies. Dash and Plotly Python documentation. https://dash.plotly.com,
        https://plotly.com/python/
  \item Waskom, M. (2021). seaborn: statistical data visualization. \emph{Journal of Open Source Software},
        6(60), 3021.
  \item DAMA International (2017). \emph{DAMA-DMBOK: Data Management Body of Knowledge}, 2nd ed.
\end{itemize}

\end{document}
